\documentclass[12pt,a4paper,openany]{book}
\usepackage[utf8]{inputenc}
\usepackage[T1,T2A]{fontenc}
\usepackage[]{graphicx} %[draft] не показывать рисы
\usepackage[russian]{babel}
\usepackage{hyperref}
\usepackage[sorting=none]{biblatex}
% \addbibresource{bibliography.bib}
\usepackage{csquotes}
\usepackage[section]{placeins}
\usepackage{float}
\usepackage[mode=buildnew]{standalone}
\usepackage[dvipsnames]{xcolor}
\definecolor{mybluelinks}{RGB}{0, 71, 211} %цвет ссылок
\hypersetup{
    colorlinks=true,
    linkcolor=mybluelinks,
    citecolor=mybluelinks,
    urlcolor=mybluelinks,
    pdftitle={Электронный справочник по физике},
    pdfauthor={Кравченко Игорь Игоревич (17. 05. 1994) (\LaTeX-код, материалы)},
    pdfpagemode=UseOutlines,%UseOutlines%UseNone
    pdfstartview=FitH,
    }

\usepackage{tikz,pgffor,bm}
\usetikzlibrary {perspective}
\usetikzlibrary {fadings,patterns,shadows}
\usetikzlibrary{positioning}
\usepackage{stackengine}
\usetikzlibrary {arrows.meta}
\usetikzlibrary {math}
\setlength\oddsidemargin{\dimexpr(\paperwidth-\textwidth)/2 - 1in\relax}
\setlength\evensidemargin{\oddsidemargin}
\usepackage[a4paper,%showframe,
            top=1in,
            bottom=1in,
            left=3cm,
            right=3cm]{geometry}
\usetikzlibrary{shadings}
\usetikzlibrary {shapes.symbols}
\usetikzlibrary {decorations.pathmorphing,decorations.pathreplacing, decorations.shapes}
\usepackage[labelsep=period,figurewithin=chapter,tablename=Табл.,tablewithin=chapter,justification=centering]{caption}[2018/05/01]
\DeclareCaptionLabelFormat{sepnum}{#1\,#2} %разделитель между номером и рис. (табл.)
\captionsetup{labelformat=sepnum}
\usepackage{amssymb,amsmath}
% \numberwithin{equation}{section}
\usetikzlibrary{intersections}
\usepackage{array}
\usepackage{scalefnt}
\usepackage[version=4]{mhchem}
\usepackage{mathtools}
\usepackage{color}
\usetikzlibrary{calc}
\usepackage{marvosym}
\usepackage{fancyhdr}
\pagestyle{fancy}
\fancyhead{}
\fancyhead[CE]{\small \href{mailto:super.antig@yandex.ru}{super.antig@yandex.ru}\hfill \href{https://vk.com/igorkravchenko1}{Кравченко Игорь Игоревич}\hfill \href{tel:+79010144910}{$+\,$7\: 901\: 014\: 49\: 10}}
\fancyhead[CO]{\small \href{mailto:pili.iveta@yandex.ru}{pili.iveta@yandex.ru}\hfill \href{https://ksdr.profi.ru/profile/KravchenkoIN3/}{Кравченко Ивета Николаевна}\hfill \href{tel:+79180260674}{$+\,$7\: 918\: 026\: 06\: 74}}
\renewcommand{\headrulewidth}{0pt}
\setlength{\headheight}{14.5pt}
%% переход с номера страницы на оглавление
\fancyfoot[C]{\hyperlink{toc}{\thepage}}
\fancypagestyle{plain}{
\fancyhead{}
\fancyfoot[C]{\hyperlink{toc}{\thepage}}
}
%%
\fancypagestyle{myfancy}{
\fancyhead[CE]{\small \href{mailto:super.antig@yandex.ru}{super.antig@yandex.ru}\hfill \href{https://vk.com/igorkravchenko1}{Кравченко Игорь Игоревич}\hfill \href{tel:+79010144910}{$+\,$7\: 901\: 014\: 49\: 10}}
\fancyhead[CO]{\small \href{mailto:pili.iveta@yandex.ru}{pili.iveta@yandex.ru}\hfill \href{https://ksdr.profi.ru/profile/KravchenkoIN3/}{Кравченко Ивета Николаевна}\hfill \href{tel:+79180260674}{$+\,$7\: 918\: 026\: 06\: 74}}
}

\usepackage{soulutf8}
\usepackage{mathrsfs}
\usepackage[europeanresistors,americaninductors,americanvoltages,compatibility]{circuitikzgit}
\usepackage[inline]{asymptote}
\usetikzlibrary{decorations.markings}
\usetikzlibrary{optics}
\usepackage{subcaption}
\usepackage{eso-pic}
\usepackage[most]{tcolorbox}
\usepackage{hhline}
\usepackage{makecell}
\renewcommand\theadfont{\normalsize}
\usepackage{enumitem}
\usepackage{wrapfig}
\usepackage{lipsum}
\usetikzlibrary{spath3}
\usetikzlibrary{angles,quotes,babel}
\usepackage{adjustbox}
\usetikzlibrary{backgrounds}
\usetikzlibrary{arrows.meta,bending}
\usepackage{tikzscale}
\usepackage{pdfpages}
\usetikzlibrary{circuits.ee.IEC}
\usepackage{pgf-spectra}


%%%%%%%%%%%%%%%%%%%%%%%%%%%%%%%%%%%%%%%
%new commands and stuff
%%%%%%%%%%%%%%%%%%%%%%%%%%%%%%%%%%%%%%%


%remove blank page of a part
\let\cleardoublepage=\clearpage

%ударение
\renewcommand{\'}{\accent1}

%перенос мат знак
\renewcommand*{\hm}[1]{#1\nobreak\discretionary{}%
{\hbox{$\mathsurround=0pt #1$}}{}}


% изменить отступ после номера параграфа в оглавлении
\makeatletter
\renewcommand*\l@section{\@dottedtocline{1}{2em}{2.3em}}
\makeatother

% изменить отступ после номера главы в оглавлении
\makeatletter
\renewcommand*\l@chapter[2]{%
  \ifnum \c@tocdepth >\m@ne
    \addpenalty{-\@highpenalty}%
    \vskip 1.0em \@plus\p@
    \setlength\@tempdima{2em}%
    \begingroup
      \parindent \z@ \rightskip \@pnumwidth
      \parfillskip -\@pnumwidth
      \leavevmode \bfseries
      \advance\leftskip\@tempdima
      \hskip -\leftskip
      #1\nobreak\hfil
      \nobreak\hb@xt@\@pnumwidth{\hss #2%
                                 \kern-\p@\kern\p@}\par
      \penalty\@highpenalty
    \endgroup
  \fi}
\makeatother


% убрать колонтитулы на страницах Часть
\makeatletter
\renewcommand\part{%
  \if@openright
    \cleardoublepage
  \else
    \clearpage
  \fi
  \thispagestyle{empty}%
  \if@twocolumn
    \onecolumn
    \@tempswatrue
  \else
    \@tempswafalse
  \fi
  \null\vfil
  \secdef\@part\@spart}
\makeatother


% убрать колонтитулы на страницах Глава
\makeatletter
\renewcommand\chapter{\if@openright\cleardoublepage\else\clearpage\fi
                    \thispagestyle{empty}%
                    \global\@topnum\z@
                    \@afterindentfalse
                    \secdef\@chapter\@schapter}
\makeatother


%boxes
\newsavebox{\myboxb}
\newlength{\mylengthb}

\newsavebox{\myboxa}
\newlength{\mylengtha}

%equation numbers
\renewcommand{\theequation}{\arabic{equation}}

%an environment law
\newenvironment{law}[1][\unskip]{%
\par
\addvspace{\medskipamount}
\noindent}
{\par\addvspace{\medskipamount}}

%pic / tab format
\renewcommand\thefigure{\arabic{figure}}
\renewcommand\thetable{\arabic{table}}

%section format
\counterwithout{section}{chapter}

% пробел после номер параграфа на странице
% \makeatletter
% \renewcommand*{\@seccntformat}[1]{\csname the#1\endcsname\hspace{1em}}
% \makeatother

%tikz установки на все картинки
\tikzset{every picture/.style={line join=round}}

%функция округления числа для tikz
\pgfmathdeclarefunction{roundn}{2}{%
\pgfmathparse{#1}%
\pgfmathparse{"\pgfmathprintnumber[fixed, precision=#2]{\pgfmathresult}"}%
}

%task comm
\newcounter{task}
\newcommand{\task}[1]{\par\addvspace{\bigskipamount}
\refstepcounter{task}
\noindent\begin{minipage}[t]{\leftmargini}
\vspace*{0pt} % вкл / выкл корректировку выравнивания minipages
\hfill\mbox{\arabic{task}.\hspace{\labelsep}}%
\end{minipage}%
\hfill\begin{minipage}[t]{\linewidth-\leftmargini}%
#1

\par\xdef\tpd{\the\prevdepth}
\end{minipage}%

\prevdepth\tpd
}


%change arrow length of {mhchem} package
\ExplSyntaxOn
\keys_define:nn { mhchem }
 {
  arrow-min-length .code:n =
   \cs_set:Npn \__mhchem_arrow_options_minLength:n { {#1} } % default is 2em
 }
\ExplSyntaxOff


% \title{Химия: атомы и химические элементы}
% \author{Кравченко Ивета Николаевна}
% \date{2021}

\addto\captionsrussian{\renewcommand*\contentsname{Оглавление}}

\usepackage{subfiles} % Best loaded last in the preamble

\begin{document}


\renewcommand\refname{Литература}

% титульные листы (варианты)
% \input{titlepage}
% \begin{titlepage}
    \begin{center}
        \vspace*{0.5cm}
        
        \Large
        Кравченко Игорь Игоревич\\
        Кравченко Ивета Николаевна
        
        \vspace{1cm}
        \Huge
        \textbf{Физика для всех:\\
        электродинамика}
        
        \vspace{1.5cm}
        
        \vfill
        \Large
        \textit{Электронный иллюстрированный справочник\\
        по физике}
        
        \vspace{15cm}

    \end{center}
\end{titlepage}
\input{title_op}


\newpage

% \addtocounter{page}{1}
\pagestyle{empty}

% убрать колонтитулы на страницах оглавления (только первая стр?)
% \addtocontents{toc}{\protect\thispagestyle{empty}} 

\phantomsection 
\hypertarget{toc}{}
\tableofcontents


\clearpage
\pagestyle{myfancy}


% \part{Механика}


% \chapter{Кинематика}


% \clearpage

% \AddToShipoutPictureFG{%
% \begin{tikzpicture}[overlay,remember picture]
% \path (current page.south west) -- (current page.north east)
% node[midway,scale=1,color=gray,sloped,opacity=0.4] {\textnormal{Кравченко Игорь Игоревич\hspace{1cm} +79010144910\hspace{1cm} super.antig@yandex.ru}};
% \end{tikzpicture}
% }


% \includepdf[pages=1,pagecommand={\addtocounter{section}{1}\phantomsection
% \addcontentsline{toc}{section}{\protect\numberline{\thesection}Векторы}}]{mech/1vec}


% \includepdf[pages=1,pagecommand={\addtocounter{section}{1}\phantomsection
% \addcontentsline{toc}{section}{\protect\numberline{\thesection}Числа и приставки}}]{mech/2num}


% \includepdf[pages=1,pagecommand={\addtocounter{section}{1}\phantomsection
% \addcontentsline{toc}{section}{\protect\numberline{\thesection}Механическое движение}}]{mech/3mot}


% \includepdf[pages=1,pagecommand={\addtocounter{section}{1}\phantomsection
% \addcontentsline{toc}{section}{\protect\numberline{\thesection}Основные понятия кинематики}}]{mech/4main}


% \includepdf[pages=1,pagecommand={\addtocounter{section}{1}\phantomsection
% \addcontentsline{toc}{section}{\protect\numberline{\thesection}Виды механического движения}}]{mech/5kinds}


% \includepdf[pages=1,pagecommand={\addtocounter{section}{1}\phantomsection
% \addcontentsline{toc}{section}{\protect\numberline{\thesection}Формулы для прямолинейного движения}}]{mech/6str}


% \includepdf[pages=1,pagecommand={\addtocounter{section}{1}\phantomsection
% \addcontentsline{toc}{section}{\protect\numberline{\thesection}Сложение и относительность скоростей}}]{mech/7add}


% \includepdf[pages=1,pagecommand={\addtocounter{section}{1}\phantomsection
% \addcontentsline{toc}{section}{\protect\numberline{\thesection}Средняя скорость}}]{mech/8ave}


% \includepdf[pages=1,pagecommand={\addtocounter{section}{1}\phantomsection
% \addcontentsline{toc}{section}{\protect\numberline{\thesection}Графики механического движения}}]{mech/9gra}


% \includepdf[pages=1,pagecommand={\addtocounter{section}{1}\phantomsection
% \addcontentsline{toc}{section}{\protect\numberline{\thesection}Свободное падение}}]{mech/10g}


% \includepdf[pages=1,pagecommand={\addtocounter{section}{1}\phantomsection
% \addcontentsline{toc}{section}{\protect\numberline{\thesection}Движение по окружности}}]{mech/11circ}


% \includepdf[pages=1,pagecommand={\addtocounter{section}{1}\phantomsection
% \addcontentsline{toc}{section}{\protect\numberline{\thesection}Основы механических передач}}]{mech/12trans}


% %%%%%%%%%%%%%%%%%%%%%%%%%%%%%%%%%%%%

% \clearpage

% \ClearShipoutPictureFG


% \chapter{Динамика}


% \clearpage

% \AddToShipoutPictureFG{%
% \begin{tikzpicture}[overlay,remember picture]
% \path (current page.south west) -- (current page.north east)
% node[midway,scale=1,color=gray,sloped,opacity=0.4] {\textnormal{Кравченко Игорь Игоревич\hspace{1cm} +79010144910\hspace{1cm} super.antig@yandex.ru}};
% \end{tikzpicture}
% }


% \includepdf[pages=1,pagecommand={\addtocounter{section}{1}\phantomsection
% \addcontentsline{toc}{section}{\protect\numberline{\thesection}Система отсчета. Действие}}]{mech/21sys}


% \includepdf[pages=1,pagecommand={\addtocounter{section}{1}\phantomsection
% \addcontentsline{toc}{section}{\protect\numberline{\thesection}Первый закон Ньютона}}]{mech/22IN}


% \includepdf[pages=1,pagecommand={\addtocounter{section}{1}\phantomsection
% \addcontentsline{toc}{section}{\protect\numberline{\thesection}Инертность и масса}}]{mech/23mass}


% \includepdf[pages=1,pagecommand={\addtocounter{section}{1}\phantomsection
% \addcontentsline{toc}{section}{\protect\numberline{\thesection}Взаимодействие и сила}}]{mech/24inter}


% \includepdf[pages=1,pagecommand={\addtocounter{section}{1}\phantomsection
% \addcontentsline{toc}{section}{\protect\numberline{\thesection}Второй и третий законы Ньютона}}]{mech/25Nlaws}


% \includepdf[pages=1,pagecommand={\addtocounter{section}{1}\phantomsection
% \addcontentsline{toc}{section}{\protect\numberline{\thesection}Силы тяготения и реакции}}]{mech/26grav}


% \includepdf[pages=1,pagecommand={\addtocounter{section}{1}\phantomsection
% \addcontentsline{toc}{section}{\protect\numberline{\thesection}Спутник}}]{mech/27sat}


% \includepdf[pages=1,pagecommand={\addtocounter{section}{1}\phantomsection
% \addcontentsline{toc}{section}{\protect\numberline{\thesection}Сила упругости}}]{mech/28elast}


% \includepdf[pages=1,pagecommand={\addtocounter{section}{1}\phantomsection
% \addcontentsline{toc}{section}{\protect\numberline{\thesection}Сила трения}}]{mech/29fric}


% \includepdf[pages=1,pagecommand={\addtocounter{section}{1}\phantomsection
% \addcontentsline{toc}{section}{\protect\numberline{\thesection}Давление}}]{mech/30press}


% %%%%%%%%%%%%%%%%%%%%%%%%%%%%%%%%%%


% \clearpage

% \ClearShipoutPictureFG


% \chapter{Статика}


% \clearpage

% \AddToShipoutPictureFG{%
% \begin{tikzpicture}[overlay,remember picture]
% \path (current page.south west) -- (current page.north east)
% node[midway,scale=1,color=gray,sloped,opacity=0.4] {\textnormal{Кравченко Игорь Игоревич\hspace{1cm} +79010144910\hspace{1cm} super.antig@yandex.ru}};
% \end{tikzpicture}
% }


% \includepdf[pages=1,pagecommand={\addtocounter{section}{1}\phantomsection
% \addcontentsline{toc}{section}{\protect\numberline{\thesection}Момент силы}}]{mech/31torq}


% \includepdf[pages=1,pagecommand={\addtocounter{section}{1}\phantomsection
% \addcontentsline{toc}{section}{\protect\numberline{\thesection}Равновесие твердого тела}}]{mech/32equil}


% \includepdf[pages=1,pagecommand={\addtocounter{section}{1}\phantomsection
% \addcontentsline{toc}{section}{\protect\numberline{\thesection}Центр масс}}]{mech/33COM}


% \includepdf[pages=1,pagecommand={\addtocounter{section}{1}\phantomsection
% \addcontentsline{toc}{section}{\protect\numberline{\thesection}Закон Паскаля}}]{mech/34Pasc}


% \includepdf[pages=1,pagecommand={\addtocounter{section}{1}\phantomsection
% \addcontentsline{toc}{section}{\protect\numberline{\thesection}Гидравлический пресс}}]{mech/35hypress}


% \includepdf[pages=1,pagecommand={\addtocounter{section}{1}\phantomsection
% \addcontentsline{toc}{section}{\protect\numberline{\thesection}Давление в жидкости}}]{mech/36liquid}


% \includepdf[pages=1,pagecommand={\addtocounter{section}{1}\phantomsection
% \addcontentsline{toc}{section}{\protect\numberline{\thesection}Закон Архимеда}}]{mech/37Arch}


% % не растягивать оставшееся на стр оглавление
% \addtocontents{toc}{\protect\clearpage}


% % % % %%%%%%%%%%%%%%%%%%%%%%%%%%%%%%%%%%


% \clearpage

% \ClearShipoutPictureFG


% \chapter{Законы сохранения}


% \clearpage

% \AddToShipoutPictureFG{%
% \begin{tikzpicture}[overlay,remember picture]
% \path (current page.south west) -- (current page.north east)
% node[midway,scale=1,color=gray,sloped,opacity=0.4] {\textnormal{Кравченко Игорь Игоревич\hspace{1cm} +79010144910\hspace{1cm} super.antig@yandex.ru}};
% \end{tikzpicture}
% }


% \includepdf[pages=1,pagecommand={\addtocounter{section}{1}\phantomsection
% \addcontentsline{toc}{section}{\protect\numberline{\thesection}Импульс}}]{mech/41imp}


% \includepdf[pages=1,pagecommand={\addtocounter{section}{1}\phantomsection
% \addcontentsline{toc}{section}{\protect\numberline{\thesection}Работа. Простые механизмы}}]{mech/42work}


% \includepdf[pages=1,pagecommand={\addtocounter{section}{1}\phantomsection
% \addcontentsline{toc}{section}{\protect\numberline{\thesection}КПД механизма}}]{mech/43ECE}


% \includepdf[pages=1,pagecommand={\addtocounter{section}{1}\phantomsection
% \addcontentsline{toc}{section}{\protect\numberline{\thesection}Энергия}}]{mech/44ener}


% \includepdf[pages=1,pagecommand={\addtocounter{section}{1}\phantomsection
% \addcontentsline{toc}{section}{\protect\numberline{\thesection}Сохранение и изменение энергии}}]{mech/45COE}


% % %%%%%%%%%%%%%%%%%%%%%%%%%%%%%%%%%%%%%


% \clearpage

% \ClearShipoutPictureFG


% \chapter{Механические колебания и волны}


% \clearpage

% \AddToShipoutPictureFG{%
% \begin{tikzpicture}[overlay,remember picture]
% \path (current page.south west) -- (current page.north east)
% node[midway,scale=1,color=gray,sloped,opacity=0.4] {\textnormal{Кравченко Игорь Игоревич\hspace{1cm} +79010144910\hspace{1cm} super.antig@yandex.ru}};
% \end{tikzpicture}
% }


% \includepdf[pages=1,pagecommand={\addtocounter{section}{1}\phantomsection
% \addcontentsline{toc}{section}{\protect\numberline{\thesection}Механические колебания}}]{mech/51oscill}


% \includepdf[pages=1,pagecommand={\addtocounter{section}{1}\phantomsection
% \addcontentsline{toc}{section}{\protect\numberline{\thesection}Маятники}}]{mech/52pend}


% \includepdf[pages=1,pagecommand={\addtocounter{section}{1}\phantomsection
% \addcontentsline{toc}{section}{\protect\numberline{\thesection}Механическая волна}}]{mech/53wave}


% % % %%%%%%%%%%%%%%%%%%%%%%%%%%%%%%%%%%%%%%%%
% % % %%%%%%%%%%%%%%%%%%%%%%%%%%%%%%%%%%%%%%%%
% % % %%%%%%%%%%%%%%%%%%%%%%%%%%%%%%%%%%%%%%%%


% \clearpage

% \ClearShipoutPictureFG


% \part{Молекулярная физика. Термодинамика}


% \chapter{Молекулярная физика}


% \clearpage

% \AddToShipoutPictureFG{%
% \begin{tikzpicture}[overlay,remember picture]
% \path (current page.south west) -- (current page.north east)
% node[midway,scale=1,color=gray,sloped,opacity=0.4] {\textnormal{Кравченко Игорь Игоревич\hspace{1cm} +79010144910\hspace{1cm} super.antig@yandex.ru}};
% \end{tikzpicture}
% }


% \includepdf[pages=1,pagecommand={\addtocounter{section}{1}\phantomsection
% \addcontentsline{toc}{section}{\protect\numberline{\thesection}Основные понятия молекулярной физики}}]{mol/1mmain}


% \includepdf[pages=1,pagecommand={\addtocounter{section}{1}\phantomsection
% \addcontentsline{toc}{section}{\protect\numberline{\thesection}Диффузия. Броуновское движение}}]{mol/2diff}


% \includepdf[pages=1,pagecommand={\addtocounter{section}{1}\phantomsection
% \addcontentsline{toc}{section}{\protect\numberline{\thesection}Основные формулы молекулярной физики}}]{mol/3mainf}


% \includepdf[pages=1,pagecommand={\addtocounter{section}{1}\phantomsection
% \addcontentsline{toc}{section}{\protect\numberline{\thesection}Формулы для идеального газа}}]{mol/4eqgas}


% \includepdf[pages=1,pagecommand={\addtocounter{section}{1}\phantomsection\addcontentsline{toc}{section}{\protect\numberline{\thesection}\texorpdfstring{Закон Д\'альтона}{Закон Да́льтона}}}]{mol/5dal}


% \includepdf[pages=1,pagecommand={\addtocounter{section}{1}\phantomsection
% \addcontentsline{toc}{section}{\protect\numberline{\thesection}Изопроцессы}}]{mol/6iso}


% \includepdf[pages=1,pagecommand={\addtocounter{section}{1}\phantomsection
% \addcontentsline{toc}{section}{\protect\numberline{\thesection}Пар}}]{mol/7vap}


% \includepdf[pages=1,pagecommand={\addtocounter{section}{1}\phantomsection
% \addcontentsline{toc}{section}{\protect\numberline{\thesection}Влажность воздуха}}]{mol/8hum}


% % % % %%%%%%%%%%%%%%%%%%%%%%%%%%%%%%%%%%%%%


% \clearpage

% \ClearShipoutPictureFG


% \chapter{Термодинамика}


% \clearpage

% \AddToShipoutPictureFG{%
% \begin{tikzpicture}[overlay,remember picture]
% \path (current page.south west) -- (current page.north east)
% node[midway,scale=1,color=gray,sloped,opacity=0.4] {\textnormal{Кравченко Игорь Игоревич\hspace{1cm} +79010144910\hspace{1cm} super.antig@yandex.ru}};
% \end{tikzpicture}
% }


% \includepdf[pages=1,pagecommand={\addtocounter{section}{1}\phantomsection
% \addcontentsline{toc}{section}{\protect\numberline{\thesection}Внутренняя энергия}}]{mol/21U}


% \includepdf[pages=1,pagecommand={\addtocounter{section}{1}\phantomsection
% \addcontentsline{toc}{section}{\protect\numberline{\thesection}Теплопередача}}]{mol/22htrans}


% \includepdf[pages=1,pagecommand={\addtocounter{section}{1}\phantomsection
% \addcontentsline{toc}{section}{\protect\numberline{\thesection}Количество теплоты}}]{mol/23Q}


% \includepdf[pages=1,pagecommand={\addtocounter{section}{1}\phantomsection
% \addcontentsline{toc}{section}{\protect\numberline{\thesection}Изменения агрегатных состояний вещества}}]{mol/24phase}


% \includepdf[pages=1,pagecommand={\addtocounter{section}{1}\phantomsection
% \addcontentsline{toc}{section}{\protect\numberline{\thesection}Графики тепловых процессов}}]{mol/25hgra}


% \includepdf[pages=1,pagecommand={\addtocounter{section}{1}\phantomsection
% \addcontentsline{toc}{section}{\protect\numberline{\thesection}Работа газа}}]{mol/26Agas}


% \includepdf[pages=1,pagecommand={\addtocounter{section}{1}\phantomsection
% \addcontentsline{toc}{section}{\protect\numberline{\thesection}Законы термодинамики}}]{mol/27LOT}


% \includepdf[pages=1,pagecommand={\addtocounter{section}{1}\phantomsection
% \addcontentsline{toc}{section}{\protect\numberline{\thesection}Тепловые машины}}]{mol/28heng}


% % % %%%%%%%%%%%%%%%%%%%%%%%%%%%%%%%%%%%%%%%%
% % % %%%%%%%%%%%%%%%%%%%%%%%%%%%%%%%%%%%%%%%%
% % % %%%%%%%%%%%%%%%%%%%%%%%%%%%%%%%%%%%%%%%%

% \clearpage

% \ClearShipoutPictureFG


% % \part{Электродинамика}


% \chapter{Электростатика}


% \clearpage

% \AddToShipoutPictureFG{%
% \begin{tikzpicture}[overlay,remember picture]
% \path (current page.south west) -- (current page.north east)
% node[midway,scale=1,color=gray,sloped,opacity=0.4] {\textnormal{Кравченко Игорь Игоревич\hspace{1cm} +79010144910\hspace{1cm} super.antig@yandex.ru}};
% \end{tikzpicture}
% }


% \includepdf[pages=1,pagecommand={\addtocounter{section}{1}\phantomsection
% \addcontentsline{toc}{section}{\protect\numberline{\thesection}Заряженные частицы и тела}}]{el/1epart}


% \includepdf[pages=1,pagecommand={\addtocounter{section}{1}\phantomsection
% \addcontentsline{toc}{section}{\protect\numberline{\thesection}Заряд}}]{el/2charge}


% \includepdf[pages=1,pagecommand={\addtocounter{section}{1}\phantomsection
% \addcontentsline{toc}{section}{\protect\numberline{\thesection}Электризация}}]{el/3electriz}


% \includepdf[pages=1,pagecommand={\addtocounter{section}{1}\phantomsection
% \addcontentsline{toc}{section}{\protect\numberline{\thesection}Электроскоп. Закон сохранения заряда}}]{el/4escope}


% \includepdf[pages=1,pagecommand={\addtocounter{section}{1}\phantomsection
% \addcontentsline{toc}{section}{\protect\numberline{\thesection}Закон Кулона. Электрическое поле}}]{el/5Claw}


% \includepdf[pages=1,pagecommand={\addtocounter{section}{1}\phantomsection
% \addcontentsline{toc}{section}{\protect\numberline{\thesection}Напряженность электрического поля}}]{el/6E}


% \includepdf[pages=1,pagecommand={\addtocounter{section}{1}\phantomsection
% \addcontentsline{toc}{section}{\protect\numberline{\thesection}Линии электрического поля}}]{el/7Eline}


% \includepdf[pages=1,pagecommand={\addtocounter{section}{1}\phantomsection
% \addcontentsline{toc}{section}{\protect\numberline{\thesection}Потенциал}}]{el/8poten}


% \includepdf[pages=1,pagecommand={\addtocounter{section}{1}\phantomsection
% \addcontentsline{toc}{section}{\protect\numberline{\thesection}Эквипотенциальные поверхности}}]{el/9eqpoten}


% \includepdf[pages=1,pagecommand={\addtocounter{section}{1}\phantomsection
% \addcontentsline{toc}{section}{\protect\numberline{\thesection}Принцип суперпозиции электрических полей}}]{el/10fprincip}


% \includepdf[pages=1,pagecommand={\addtocounter{section}{1}\phantomsection
% \addcontentsline{toc}{section}{\protect\numberline{\thesection}Проводник}}]{el/11cond}


% \includepdf[pages=1,pagecommand={\addtocounter{section}{1}\phantomsection
% \addcontentsline{toc}{section}{\protect\numberline{\thesection}Проводник в электрическом поле}}]{el/12condef}


% \includepdf[pages=1,pagecommand={\addtocounter{section}{1}\phantomsection
% \addcontentsline{toc}{section}{\protect\numberline{\thesection}Электрическое поле проводящего шара}}]{el/13esphere}


% \includepdf[pages=1,pagecommand={\addtocounter{section}{1}\phantomsection
% \addcontentsline{toc}{section}{\protect\numberline{\thesection}Диэлектрик}}]{el/14diel}


% \includepdf[pages=1,pagecommand={\addtocounter{section}{1}\phantomsection
% \addcontentsline{toc}{section}{\protect\numberline{\thesection}Конденсатор}}]{el/15cap}


% \includepdf[pages=1,pagecommand={\addtocounter{section}{1}\phantomsection
% \addcontentsline{toc}{section}{\protect\numberline{\thesection}Соединения конденсаторов}}]{el/16caps}


% % % %%%%%%%%%%%%%%%%%%%%%%%%%%%%%%%%%%%%%


% \clearpage

% \ClearShipoutPictureFG


% \chapter{Постоянный ток}


% \clearpage

% \AddToShipoutPictureFG{%
% \begin{tikzpicture}[overlay,remember picture]
% \path (current page.south west) -- (current page.north east)
% node[midway,scale=1,color=gray,sloped,opacity=0.4] {\textnormal{Кравченко Игорь Игоревич\hspace{1cm} +79010144910\hspace{1cm} super.antig@yandex.ru}};
% \end{tikzpicture}
% }


% \includepdf[pages=1,pagecommand={\addtocounter{section}{1}\phantomsection
% \addcontentsline{toc}{section}{\protect\numberline{\thesection}Электрический ток}}]{el/21I}


% \includepdf[pages=1,pagecommand={\addtocounter{section}{1}\phantomsection
% \addcontentsline{toc}{section}{\protect\numberline{\thesection}Сопротивление}}]{el/22R}


% \includepdf[pages=1,pagecommand={\addtocounter{section}{1}\phantomsection
% \addcontentsline{toc}{section}{\protect\numberline{\thesection}Закон Ома}}]{el/23Ohm}


% \includepdf[pages=1,pagecommand={\addtocounter{section}{1}\phantomsection
% \addcontentsline{toc}{section}{\protect\numberline{\thesection}Соединения проводников}}]{el/24conn}


% \includepdf[pages=1,pagecommand={\addtocounter{section}{1}\phantomsection
% \addcontentsline{toc}{section}{\protect\numberline{\thesection}Работа тока. Закон Джоуля"--~Ленца}}]{el/25Acurr}


% \includepdf[pages=1,pagecommand={\addtocounter{section}{1}\phantomsection
% \addcontentsline{toc}{section}{\protect\numberline{\thesection}Источник тока}}]{el/26Eforce}


% \includepdf[pages=1,pagecommand={\addtocounter{section}{1}\phantomsection
% \addcontentsline{toc}{section}{\protect\numberline{\thesection}Закон Ома для полной цепи}}]{el/27fOhm}


% \includepdf[pages=1,pagecommand={\addtocounter{section}{1}\phantomsection
% \addcontentsline{toc}{section}{\protect\numberline{\thesection}Полупроводник}}]{el/28semi}


% \includepdf[pages=1,pagecommand={\addtocounter{section}{1}\phantomsection
% \addcontentsline{toc}{section}{\protect\numberline{\thesection}Полупроводники \emph{p}- и \emph{n}-типа}}]{el/28semi2}


% \includepdf[pages=1,pagecommand={\addtocounter{section}{1}\phantomsection
% \addcontentsline{toc}{section}{\protect\numberline{\thesection}\emph{p}--\emph{n}-Переход}}]{el/29pn}


% % % %%%%%%%%%%%%%%%%%%%%%%%%%%%%%%%%%%%%%


% \clearpage

% \ClearShipoutPictureFG


% \chapter{Магнитное поле}


% \clearpage

% \AddToShipoutPictureFG{%
% \begin{tikzpicture}[overlay,remember picture]
% \path (current page.south west) -- (current page.north east)
% node[midway,scale=1,color=gray,sloped,opacity=0.4] {\textnormal{Кравченко Игорь Игоревич\hspace{1cm} +79010144910\hspace{1cm} super.antig@yandex.ru}};
% \end{tikzpicture}
% }


% \includepdf[pages=1,pagecommand={\addtocounter{section}{1}\phantomsection
% \addcontentsline{toc}{section}{\protect\numberline{\thesection}Магнит. Магнитное поле}}]{el/31magn}


% \includepdf[pages=1,pagecommand={\addtocounter{section}{1}\phantomsection
% \addcontentsline{toc}{section}{\protect\numberline{\thesection}Магнитная стрелка}}]{el/32magnne}


% \includepdf[pages=1,pagecommand={\addtocounter{section}{1}\phantomsection
% \addcontentsline{toc}{section}{\protect\numberline{\thesection}Индукция и линии магнитного поля}}]{el/33B}


% \includepdf[pages=1,pagecommand={\addtocounter{section}{1}\phantomsection
% \addcontentsline{toc}{section}{\protect\numberline{\thesection}Принцип суперпозиции магнитных полей}}]{el/34Bprinc}


% \includepdf[pages=1,pagecommand={\addtocounter{section}{1}\phantomsection
% \addcontentsline{toc}{section}{\protect\numberline{\thesection}Опыт Эрстеда}}]{el/35Erst}


% \includepdf[pages=1,pagecommand={\addtocounter{section}{1}\phantomsection
% \addcontentsline{toc}{section}{\protect\numberline{\thesection}Магнитное поле прямого провода с током}}]{el/36wire}


% \includepdf[pages=1,pagecommand={\addtocounter{section}{1}\phantomsection
% \addcontentsline{toc}{section}{\protect\numberline{\thesection}Магнитное поле кольца с током}}]{el/37ring}


% % не растягивать оставшееся на стр оглавление
% \addtocontents{toc}{\protect\clearpage}


% \includepdf[pages=1,pagecommand={\addtocounter{section}{1}\phantomsection
% \addcontentsline{toc}{section}{\protect\numberline{\thesection}Магнитное поле катушки с током}}]{el/38coil}


% \includepdf[pages=1,pagecommand={\addtocounter{section}{1}\phantomsection
% \addcontentsline{toc}{section}{\protect\numberline{\thesection}Сила Ампера}}]{el/39aFA}


% \includepdf[pages=1,pagecommand={\addtocounter{section}{1}\phantomsection
% \addcontentsline{toc}{section}{\protect\numberline{\thesection}Сила Лоренца}}]{el/39bFL}


% % % %%%%%%%%%%%%%%%%%%%%%%%%%%%%%%%%%%%%%


% \clearpage

% \ClearShipoutPictureFG


% \chapter{Электромагнитная индукция}


% \clearpage

% \AddToShipoutPictureFG{%
% \begin{tikzpicture}[overlay,remember picture]
% \path (current page.south west) -- (current page.north east)
% node[midway,scale=1,color=gray,sloped,opacity=0.4] {\textnormal{Кравченко Игорь Игоревич\hspace{1cm} +79010144910\hspace{1cm} super.antig@yandex.ru}};
% \end{tikzpicture}
% }


% \includepdf[pages=1,pagecommand={\addtocounter{section}{1}\phantomsection
% \addcontentsline{toc}{section}{\protect\numberline{\thesection}Магнитный поток}}]{el/41mflux}


% \includepdf[pages=1,pagecommand={\addtocounter{section}{1}\phantomsection
% \addcontentsline{toc}{section}{\protect\numberline{\thesection}Закон электромагнитной индукции Фарадея}}]{el/42Flaw}


% \includepdf[pages=1,pagecommand={\addtocounter{section}{1}\phantomsection
% \addcontentsline{toc}{section}{\protect\numberline{\thesection}ЭДС индукции в движущемся проводе}}]{el/43mcond}


% \includepdf[pages=1,pagecommand={\addtocounter{section}{1}\phantomsection
% \addcontentsline{toc}{section}{\protect\numberline{\thesection}Правило Ленца}}]{el/44Lenz}


% \includepdf[pages=1,pagecommand={\addtocounter{section}{1}\phantomsection
% \addcontentsline{toc}{section}{\protect\numberline{\thesection}Индуктивность. Самоиндукция}}]{el/45si}


% % %%%%%%%%%%%%%%%%%%%%%%%%%%%%%%%%%%%%%


% \clearpage

% \ClearShipoutPictureFG


% \chapter[\texorpdfstring{Электромагнитные колебания и волны}{Электромагнитные колебания и волны}]{Электромагнитные колебания\\ и волны}


% \clearpage

% \AddToShipoutPictureFG{%
% \begin{tikzpicture}[overlay,remember picture]
% \path (current page.south west) -- (current page.north east)
% node[midway,scale=1,color=gray,sloped,opacity=0.4] {\textnormal{Кравченко Игорь Игоревич\hspace{1cm} +79010144910\hspace{1cm} super.antig@yandex.ru}};
% \end{tikzpicture}
% }



% \includepdf[pages=1,pagecommand={\addtocounter{section}{1}\phantomsection
% \addcontentsline{toc}{section}{\protect\numberline{\thesection}Электромагнитные колебания}}]{el/51emos}


% \includepdf[pages=1,pagecommand={\addtocounter{section}{1}\phantomsection
% \addcontentsline{toc}{section}{\protect\numberline{\thesection}Свободные колебания в контуре}}]{el/52LC}


% \includepdf[pages=1,pagecommand={\addtocounter{section}{1}\phantomsection
% \addcontentsline{toc}{section}{\protect\numberline{\thesection}Вынужденные колебания в контуре}}]{el/53dros}


% \includepdf[pages=1,pagecommand={\addtocounter{section}{1}\phantomsection
% \addcontentsline{toc}{section}{\protect\numberline{\thesection}Электроэнергия}}]{el/54EE}


% \includepdf[pages=1,pagecommand={\addtocounter{section}{1}\phantomsection
% \addcontentsline{toc}{section}{\protect\numberline{\thesection}Электромагнитная волна}}]{el/55ewave}


% % %%%%%%%%%%%%%%%%%%%%%%%%%%%%%%%%%%%%%


\clearpage

\ClearShipoutPictureFG


\part{Оптика}


\chapter{Геометрическая оптика}


\clearpage

\AddToShipoutPictureFG{%
\begin{tikzpicture}[overlay,remember picture]
\path (current page.south west) -- (current page.north east)
node[midway,scale=1,color=gray,sloped,opacity=0.4] {\textnormal{Кравченко Игорь Игоревич\hspace{1cm} +79010144910\hspace{1cm} super.antig@yandex.ru}};
\end{tikzpicture}
}



\includepdf[pages=1,pagecommand={\addtocounter{section}{1}\phantomsection
\addcontentsline{toc}{section}{\protect\numberline{\thesection}Свет. Луч света}}]{op/1light}


\includepdf[pages=1,pagecommand={\addtocounter{section}{1}\phantomsection
\addcontentsline{toc}{section}{\protect\numberline{\thesection}Отражение света. Зеркало}}]{op/2refl}


\includepdf[pages=1,pagecommand={\addtocounter{section}{1}\phantomsection
\addcontentsline{toc}{section}{\protect\numberline{\thesection}Преломление света}}]{op/3refr}


\includepdf[pages=1,pagecommand={\addtocounter{section}{1}\phantomsection
\addcontentsline{toc}{section}{\protect\numberline{\thesection}Полное внутреннее отражение}}]{op/4tref}


\includepdf[pages=1,pagecommand={\addtocounter{section}{1}\phantomsection
\addcontentsline{toc}{section}{\protect\numberline{\thesection}Линза}}]{op/5lens}


\includepdf[pages=1,pagecommand={\addtocounter{section}{1}\phantomsection
\addcontentsline{toc}{section}{\protect\numberline{\thesection}Ход лучей в тонкой линзе}}]{op/6tlens}


\includepdf[pages=1,pagecommand={\addtocounter{section}{1}\phantomsection
\addcontentsline{toc}{section}{\protect\numberline{\thesection}Построение изображения в линзе}}]{op/7image}


\includepdf[pages=1,pagecommand={\addtocounter{section}{1}\phantomsection
\addcontentsline{toc}{section}{\protect\numberline{\thesection}Глаз}}]{op/8eye}


%%%%%%%%%%%%%%%%%%%%%%%


\clearpage

\ClearShipoutPictureFG


\chapter{Волновая оптика}


\clearpage

\AddToShipoutPictureFG{%
\begin{tikzpicture}[overlay,remember picture]
\path (current page.south west) -- (current page.north east)
node[midway,scale=1,color=gray,sloped,opacity=0.4] {\textnormal{Кравченко Игорь Игоревич\hspace{1cm} +79010144910\hspace{1cm} super.antig@yandex.ru}};
\end{tikzpicture}
}


\includepdf[pages=1,pagecommand={\addtocounter{section}{1}\phantomsection
\addcontentsline{toc}{section}{\protect\numberline{\thesection}Сложение световых волн}}]{op/21lprinc}


\includepdf[pages=1,pagecommand={\addtocounter{section}{1}\phantomsection
\addcontentsline{toc}{section}{\protect\numberline{\thesection}Интерференция света}}]{op/22linter}


\includepdf[pages=1,pagecommand={\addtocounter{section}{1}\phantomsection
\addcontentsline{toc}{section}{\protect\numberline{\thesection}Дифракция света}}]{op/23diffr}


\includepdf[pages=1,pagecommand={\addtocounter{section}{1}\phantomsection
\addcontentsline{toc}{section}{\protect\numberline{\thesection}Дисперсия света}}]{op/24disp}


% % %%%%%%%%%%%%%%%%%%%%%%%%%%%%%%%%%%%%%%%%
% % %%%%%%%%%%%%%%%%%%%%%%%%%%%%%%%%%%%%%%%%
% % %%%%%%%%%%%%%%%%%%%%%%%%%%%%%%%%%%%%%%%%


\clearpage

\ClearShipoutPictureFG


\part{Теория относительности}


\chapter[Основы специальной теории относительности (СТО)]{Основы специальной\\ теории относительности (СТО)}


\clearpage

\AddToShipoutPictureFG{%
\begin{tikzpicture}[overlay,remember picture]
\path (current page.south west) -- (current page.north east)
node[midway,scale=1,color=gray,sloped,opacity=0.4] {\textnormal{Кравченко Игорь Игоревич\hspace{1cm} +79010144910\hspace{1cm} super.antig@yandex.ru}};
\end{tikzpicture}
}


\includepdf[pages=1,pagecommand={\addtocounter{section}{1}\phantomsection
\addcontentsline{toc}{section}{\protect\numberline{\thesection}Принципы СТО}}]{sr/1sprinc}


\includepdf[pages=1,pagecommand={\addtocounter{section}{1}\phantomsection
\addcontentsline{toc}{section}{\protect\numberline{\thesection}Основные формулы СТО}}]{sr/2smain}


% % %%%%%%%%%%%%%%%%%%%%%%%%%%%%%%%%%%%%%%%%
% % %%%%%%%%%%%%%%%%%%%%%%%%%%%%%%%%%%%%%%%%
% % %%%%%%%%%%%%%%%%%%%%%%%%%%%%%%%%%%%%%%%%


\clearpage

\ClearShipoutPictureFG


\part{Квантовая физика}


\chapter{Фотоны}


\clearpage

\AddToShipoutPictureFG{%
\begin{tikzpicture}[overlay,remember picture]
\path (current page.south west) -- (current page.north east)
node[midway,scale=1,color=gray,sloped,opacity=0.4] {\textnormal{Кравченко Игорь Игоревич\hspace{1cm} +79010144910\hspace{1cm} super.antig@yandex.ru}};
\end{tikzpicture}
}



\includepdf[pages=1,pagecommand={\addtocounter{section}{1}\phantomsection
\addcontentsline{toc}{section}{\protect\numberline{\thesection}Фотоэффект}}]{qu/1photo}


\includepdf[pages=1,pagecommand={\addtocounter{section}{1}\phantomsection
\addcontentsline{toc}{section}{\protect\numberline{\thesection}Формулы для фотоэффекта}}]{qu/2pheq}


\includepdf[pages=1,pagecommand={\addtocounter{section}{1}\phantomsection
\addcontentsline{toc}{section}{\protect\numberline{\thesection}Давление света}}]{qu/3lpress}


% % %%%%%%%%%%%%%%%%%%%%%%%%%%%%%%%%%%%%%


\clearpage

\ClearShipoutPictureFG


\chapter{Физика атома}


\clearpage

\AddToShipoutPictureFG{%
\begin{tikzpicture}[overlay,remember picture]
\path (current page.south west) -- (current page.north east)
node[midway,scale=1,color=gray,sloped,opacity=0.4] {\textnormal{Кравченко Игорь Игоревич\hspace{1cm} +79010144910\hspace{1cm} super.antig@yandex.ru}};
\end{tikzpicture}
}


\includepdf[pages=1,pagecommand={\addtocounter{section}{1}\phantomsection
\addcontentsline{toc}{section}{\protect\numberline{\thesection}Постулаты Бора}}]{qu/4bohr}


\includepdf[pages=1,pagecommand={\addtocounter{section}{1}\phantomsection
\addcontentsline{toc}{section}{\protect\numberline{\thesection}Спектры}}]{qu/5spec}


% % %%%%%%%%%%%%%%%%%%%%%%%%%%%%%%%%%%%%%


\clearpage

\ClearShipoutPictureFG


\chapter{Физика атомного ядра}


\clearpage

\AddToShipoutPictureFG{%
\begin{tikzpicture}[overlay,remember picture]
\path (current page.south west) -- (current page.north east)
node[midway,scale=1,color=gray,sloped,opacity=0.4] {\textnormal{Кравченко Игорь Игоревич\hspace{1cm} +79010144910\hspace{1cm} super.antig@yandex.ru}};
\end{tikzpicture}
}



\includepdf[pages=1,pagecommand={\addtocounter{section}{1}\phantomsection
\addcontentsline{toc}{section}{\protect\numberline{\thesection}Атомное ядро}}]{qu/6nucl}


\includepdf[pages=1,pagecommand={\addtocounter{section}{1}\phantomsection
\addcontentsline{toc}{section}{\protect\numberline{\thesection}Радиоактивность. Ядерные реакции}}]{qu/7rad}


% % не растягивать оставшееся на стр оглавление
\addtocontents{toc}{\protect\clearpage}


% % %%%%%%%%%%%%%%%%%%%%%%%%%%%%%%%%%%%%%%%%
% % %%%%%%%%%%%%%%%%%%%%%%%%%%%%%%%%%%%%%%%%
% % %%%%%%%%%%%%%%%%%%%%%%%%%%%%%%%%%%%%%%%%


% \part{Астрофизика}


% \chapter{Элементы астрофизики}


% \newpage
% \section{Солнечная система}


% \newpage
% \section{Звезда}

% \newpage
% \section{Звездные характеристики}


% \newpage
% \section{Закономерности звездных характеристик}


% \newpage
% \section{О происхождении и эволюции звезд}


% \newpage
% \section{Схема эволюции звезд}


% \newpage
% \section{Масштабы и эволюция Вселенной}


% \newpage
% \section{Строение Вселенной}


% % %%%%%%%%%%%%%%%%%%%%%%%%%%%%%%%%%%%%%%%%
% % %%%%%%%%%%%%%%%%%%%%%%%%%%%%%%%%%%%%%%%%
% % %%%%%%%%%%%%%%%%%%%%%%%%%%%%%%%%%%%%%%%%


\clearpage

\ClearShipoutPictureFG


\part{Приложение}


\chapter{Как решать задачи по физике}


\clearpage

\AddToShipoutPictureFG{%
\begin{tikzpicture}[overlay,remember picture]
\path (current page.south west) -- (current page.north east)
node[midway,scale=1,color=gray,sloped,opacity=0.4] {\textnormal{Кравченко Игорь Игоревич\hspace{1cm} +79010144910\hspace{1cm} super.antig@yandex.ru}};
\end{tikzpicture}
}


\includepdf[pages=1,pagecommand={\addtocounter{section}{1}\phantomsection
\addcontentsline{toc}{section}{\protect\numberline{\thesection}Как находить ответ к задаче по физике}}]{za/find}


\includepdf[pages=1,pagecommand={\addtocounter{section}{1}\phantomsection
\addcontentsline{toc}{section}{\protect\numberline{\thesection}Решение специфических задач по физике}}]{za/mnote}


% % %%%%%%%%%%%%%%%%%%%%%%%%%%%%%%%%%%%%%%%%
% \chapter{Приборы в физике}


% \newpage
% \section{Приборы. Механические явления}


% \newpage
% \section{Приборы. Тепловые явления}


% \newpage
% \section{Приборы. Электромагнитные явления}


% % %%%%%%%%%%%%%%%%%%%%%%%%%%%%%%%%%%%%%%%%
% \chapter{Элементы математики}


% \newpage
% \section{Углы в математике}


% \newpage
% \section{Площади фигур в математике}


% \newpage
% \section{Объемы тел в математике}


% \newpage
% \section{Производная}


% % %%%%%%%%%%%%%%%%%%%%%%%%%%%%%%%%%%%%%%%%
% \chapter{Астрофика. Справочный материал}


% \newpage
% \section{Основные величины в астрофизике}














\newpage
\thispagestyle{plain}

% \printbibliography

% \addcontentsline{toc}{section}{Литература}

\end{document}
